% 第4章 実験結果
% 対応する草稿: research/thesis-drafts/chapter4-results-skeleton.md

\subsection{実験設定のまとめ}\label{subsec:setup-summary}

% TODO: データ取得後に記述

\subsection{全体傾向（実験1）}\label{subsec:overall}

% TODO: モデル別 calibration metrics 表
% Reliability Diagram (Figure)

\subsection{ドメイン別の分析}\label{subsec:domain-analysis}

% TODO: ドメイン × モデルの ECE 棒グラフ
% 興味深い事例

\subsection{プロンプト方式の影響（実験2）}\label{subsec:prompt-analysis}

% TODO: プロンプト × モデルの ECE 棒グラフ

\subsection{日英言語の比較（実験3）}\label{subsec:language-comparison}

% TODO: 日本語 vs 英語の Reliability Diagram

\subsection{追加分析}\label{subsec:additional}

% TODO: 信頼度分布ヒストグラム、誤答時信頼度、AUROC

\subsection{本章のまとめ}\label{subsec:results-summary}

% TODO: 主要発見3点の要約
